% 李群 SO(3) 示意 (Lie Group SO(3))
% 描述: 展示三维旋转群的多种表示
% 数学概念: 特殊正交群、旋转、欧拉角、四元数、Bloch球
\documentclass[border=10pt]{standalone}
\usepackage{tikz}
\usepackage{amsmath,amssymb}
\usetikzlibrary{arrows.meta,positioning,calc,3d,patterns,decorations.pathmorphing}

\begin{document}
\begin{tikzpicture}[
    >=Stealth,
    scale=1.0
]

% ============= 左侧: 旋转作用 =============
\begin{scope}[xshift=-5cm, yshift=1cm]
    % 标题
    \node[font=\bfseries] at (0, 3) {旋转变换};
    
    % 旋转前
    \draw[color=gray!50, line width=1pt] (-1.5, 1) -- (1.5, 1);
    \draw[color=gray!50, line width=1pt] (0, 0) -- (0, 2);
    \draw[color=gray!50, line width=1pt] (-0.8, 0.2) -- (0.8, 1.8);
    \node[color=gray!50, font=\tiny] at (1.7, 1) {原位置};
    
    % 旋转后
    \begin{scope}[rotate around={45:(0, 1)}]
        \draw[color=blue!60!black, line width=1.5pt] (-1.5, 1) -- (1.5, 1);
        \draw[color=green!60!black, line width=1.5pt] (0, 0) -- (0, 2);
        \draw[color=red!70!black, line width=1.5pt] (-0.8, 0.2) -- (0.8, 1.8);
    \end{scope}
    
    % 旋转中心
    \fill[black] (0, 1) circle (2pt);
    
    % 旋转箭头
    \draw[->, color=orange!80!black, line width=1.5pt] (0.8, 0.3) arc (-45:45:0.8);
    \node[color=orange!80!black, font=\small, right] at (1.2, 0.5) {$R$};
    
    % 矩阵表示
    \node[anchor=north, font=\scriptsize] at (0, -0.5) {$R \in SO(3)$};
    \node[anchor=north, font=\tiny, text width=4cm, align=center] 
        at (0, -1.2) {$R^T R = I, \; \det R = 1$};
\end{scope}

% ============= 中间: 欧拉角 =============
\begin{scope}[xshift=0cm, yshift=1cm]
    % 标题
    \node[font=\bfseries] at (0, 3) {欧拉角分解};
    
    % 坐标系 (初始)
    \draw[->, color=gray!50, line width=1pt] (-1.5, -1) -- (-0.5, -1) node[right, font=\tiny] {$x$};
    \draw[->, color=gray!50, line width=1pt] (-1.5, -1) -- (-1.5, 0) node[above, font=\tiny] {$z$};
    
    % 旋转步骤
    \node[anchor=north west, font=\scriptsize, text width=4cm, align=left] 
        at (-1.8, -0.3) {
        1. 绕 $z$ 轴旋转 $\phi$ \\
        2. 绕新 $x'$ 轴旋转 $\theta$ \\
        3. 绕新 $z''$ 轴旋转 $\psi$
    };
    
    % 最终坐标系 (旋转后)
    \begin{scope}[rotate around={30:(1, 0.5)}]
        \draw[->, color=blue!60!black, line width=1pt] (1, 0.5) -- (2, 0.5) node[right, font=\tiny] {$x'''$};
        \draw[->, color=green!60!black, line width=1pt] (1, 0.5) -- (1, 1.5) node[above, font=\tiny] {$z'''$};
    \end{scope}
    
    % 公式
    \node[anchor=north, font=\scriptsize] at (0, -2) {$R = R_z(\psi) R_x(\theta) R_z(\phi)$};
    
    % 说明
    \node[anchor=north, font=\tiny, text width=4.5cm, align=center] 
        at (0, -2.8) {$\phi, \psi \in [0, 2\pi), \theta \in [0, \pi]$};
\end{scope}

% ============= 右侧: 四元数表示 =============
\begin{scope}[xshift=5cm, yshift=1cm]
    % 标题
    \node[font=\bfseries] at (0, 3) {四元数表示};
    
    % 球面 (S³)
    \draw[color=blue!60!black, line width=1pt, fill=blue!10, opacity=0.3] 
        (0, 0.5) circle (1.5);
    
    % 赤道圆
    \draw[color=green!60!black, line width=0.8pt] 
        (0, 0.5) ellipse (1.5 and 0.4);
    
    % 四元数点
    \coordinate (Q) at (0.8, 1.2);
    \fill[red!70!black] (Q) circle (3pt);
    \node[color=red!70!black, font=\small, above right] at (Q) {$q$};
    
    % 对径点 (等同)
    \coordinate (Qneg) at (-0.8, -0.2);
    \fill[red!70!black] (Qneg) circle (3pt);
    \node[color=red!70!black, font=\small, below left] at (Qneg) {$-q$};
    
    % 等同标记
    \draw[<->, color=orange!80!black, line width=1pt] (0.5, 0.9) -- (-0.5, 0.1);
    \node[color=orange!80!black, font=\tiny] at (0, 0.9) {等同};
    
    % 公式
    \node[anchor=north, font=\scriptsize] at (0, -1.2) {$q = a + bi + cj + dk$};
    \node[anchor=north, font=\scriptsize] at (0, -1.8) {$|q| = 1$};
    
    % 关系
    \node[anchor=north, font=\tiny, text width=4.5cm, align=center] 
        at (0, -2.8) {$SO(3) \cong S^3 / \{q \sim -q\} \cong \mathbb{R}P^3$};
\end{scope}

% ============= 底部: 李代数结构 =============
\node[anchor=north, font=\small, draw=blue!30, fill=blue!5, rounded corners, inner sep=10pt] 
    at (4, -4) {
    \begin{tabular}{ll}
        \textbf{SO(3) 的群结构:} & \\
        流形结构: & 三维实射影空间 $\mathbb{R}P^3$ \\
        维数: & $\dim SO(3) = 3$ \\
        李代数: & $\mathfrak{so}(3) \cong (\mathbb{R}^3, \times)$ (叉积) \\
        指数映射: & $\exp: \mathfrak{so}(3) \to SO(3)$ \\
        单连通覆叠: & $SU(2) \to SO(3)$ (2:1覆叠) \\
        基本群: & $\pi_1(SO(3)) \cong \mathbb{Z}_2$
    \end{tabular}
};

% ============= 无穷小生成元 =============
\begin{scope}[xshift=-5cm, yshift=-4cm]
    \node[anchor=north west, font=\scriptsize, text width=5cm, align=left] 
        at (0, 0) {
        \textbf{无穷小生成元:}\\[5pt]
        $J_x = \begin{pmatrix} 0 & 0 & 0 \\ 0 & 0 & -1 \\ 0 & 1 & 0 \end{pmatrix}$\\[8pt]
        $J_y = \begin{pmatrix} 0 & 0 & 1 \\ 0 & 0 & 0 \\ -1 & 0 & 0 \end{pmatrix}$\\[8pt]
        $J_z = \begin{pmatrix} 0 & -1 & 0 \\ 1 & 0 & 0 \\ 0 & 0 & 0 \end{pmatrix}$
    };
\end{scope}

% ============= 物理应用 =============
\begin{scope}[xshift=5cm, yshift=-4cm]
    \node[anchor=north east, font=\scriptsize, text width=5cm, align=left] 
        at (0, 0) {
        \textbf{物理应用:}\\
        - 刚体转动 \\
        - 角动量量子化 \\
        - 自旋1/2粒子 \\
        - 陀螺仪导航 \\
        - 计算机图形学\\[5pt]
        \textit{Dirac鞋带把戏}\\
        证明$\pi_1(SO(3)) = \mathbb{Z}_2$
    };
\end{scope}

\end{tikzpicture}
\end{document}
